\section*{Integral Aufgaben mit L�sungen}\label{MA2-integration-loesungen}

\ifintegriereaufgaben Erkl�rungen und Beispiele, die zum Verst�ndis der L�sungswege hilfreich sind, finden Sie im Skript\wennreferenzvorhanden{MA2-integration}{ (Abschnitt \abschnittsverweis{MA2-integration})} des Skripts.\fi

\begin{exercisebox}[Fl�cheninhalt]
	Berechnen mit Hilfe geometrischer �berlegungen den Fl�cheninhalt zwischen Graphen und $x$-Achse

	\begin{eqnarray*}
		(a) \;  \displaystyle{ \int\limits_0^4} (2x-1) \mathrm{d}x \quad  \quad  (b) \; \displaystyle{ \int\limits_0^2} (x-1)^3 \mathrm{d}x  \quad \quad (c) \; \displaystyle{ \int\limits_0^1} |\sin(2\pi x)| \mathrm{d}x
	\end{eqnarray*}

\vspace{0.3cm}
\textbf{L�sung:} \newline
\begin{itemize}
	\item[(a)] $f: x \mapsto 2x-1, \quad f(\frac12) = 0 \Rightarrow A = -\frac12\cdot \frac12 \cdot 1 + \frac12 \cdot 7 \cdot (4-\frac12) = 12 \,. $
	\item[(b)] $\forall x \in [0,1]: \quad f(x) = -f(x+1) \Rightarrow A = 0$
	\item[(c)] $ A = \frac12 1 \cdot 1 = \frac12$
\end{itemize}
\end{exercisebox}

\begin{exercisebox}[Fl�cheninhalt (Trapezformel)] % LS p 153
	\begin{itemize}
		\item[(a)] $A$ sei der Fl�cheninhalt eines Kreises mit Radius $r$. Zeigen Sie: $ 2r^2 < A < 4 r^2$.
		\item[(b)] �berlegen Sie sich, wie Sie den Fl�chenhinhalt $A$ n�herhungsweise mit einer numerischen Integration berechnen k�nnen.
		\item[(c)] Was ist der Unterschied zur Fl�cheninhaltsberechnung von Vielecken?
	\end{itemize}

\vspace{0.3cm}
\textbf{L�sung:} \newline

\begin{itemize}
	\item[(a)] Einbeschriebenes Quadrat hat die Fl�che \( (\sqrt{2}r)^2 \) und das umschreibende Quadrat die Fl�che \( (\sqrt{4} r)^2\).
	\item[(b)] Eine Idee ist die Fl�che unter einem Halbkreis zu berechnen.  Der Halbkreis ist n�mlich als Funktion darstellbar und durch \(f: x\mapsto \sqrt{r^2-x^2}\) zu parametrisieren. Den Fl�cheninhalt kann man durch eine Zerlegungssumme approximieren, beispielsweise durch die Trapezformel. Dazu unterteilen wir das Intervall \( [-r,r]\) in \(n\in\mathbb N\) Teilintervalle mit Schrittweite $h=2r/n$ also \(x_k = -r + k\cdot h\) und berechnen sie Zerlegungssumme \(S_n\):\footnote{Sp�ter lernen wir die numerische Integration mit Zufallszahlen kennen, die sogenannte Monte-Carlo Methode}
	\[ S_n = \sum\limits_{k=0}^{n}  \frac{ f( x_{k+1}) + f(x_k)}{2} \cdot h  \Rightarrow  A \approx \lim\limits_{n\to \infty} 2\cdot S_n = \pi r^2\,.\]
\item[(c)] Anders als bei Vielecken, kann der Fl�cheninhalt nicht exakt bestimmt werden (die Kr�mmung des Funktionsgraphen ist nie konstant!). Um die Fl�chenberechnung zu berechnen muss ein Grenzwert gebildet werden.
	\end{itemize}
\end{exercisebox}
\begin{exercisebox}[Eigenschaften des Integrals]
	Begr�nden Sie, dass jede Integralfunktion $ F(x):= \displaystyle{\int\limits_{a}^x} f(t)\,\mathrm{d}t$ mindestens eine Nullstelle hat.

\vspace{0.3cm}
\textbf{L�sung:} Das ist eine Anwendung der Eigenschaften des Riemann-Integrals, denn es gilt $ \displaystyle{\int\limits_{a}^a} f(t)\,\mathrm{d}t = 0$.
\end{exercisebox}


\begin{exercisebox}[Hauptsatz der Differential- und Integralrechnung]  %Papula 1, V, Abschnitt 1-7, Aufg. 1
	Bestimmen Sie alle Stammfunktionen folgender Funktionen:
	\begin{eqnarray*}
		\begin{array}{l l l }
			(a) \;  x\mapsto 4x^5-6x^3+8x^2-3x +5
			& (b) \; x\mapsto e^x + x^2 - 2x + \sin(x)
			& (c)\; x\mapsto \displaystyle{\frac{\tan(x)}{\sin(2x)}}
		\end{array}
	\end{eqnarray*}

\vspace{0.3cm}
\textbf{L�sung:} \newline
Die hier gestellten Aufgaben (und viele mehr :) ) finden Sie in Papula, Band 1, Abschnitt V.1-7. Aufgaben ab Seite \(559\) und die L�sungen auf Seite \(787\)
\begin{itemize}
\item[(a)] Papula Aufg. 1.a) $\quad F(x) = \displaystyle{\frac23} x^6 - \displaystyle{\frac32}x^4 + \displaystyle{\frac83} x^3 - \displaystyle{\frac32} x^2 + 5x + C $
	\item[(b)] Papula Aufg. 2.a) $\quad F(x) = e^x + \displaystyle{\frac13} x^3 - x^2 - \cos(x) + C $
	\item[(c)] Papula Aufg. 2.j) $\quad F(x) = \displaystyle{\frac12}\tan(x) + C $
\end{itemize}
\end{exercisebox}

\begin{exercisebox}[Bestimmte und uneigentliche Integrale] %Papula 1, V, Abschnitt 1-7, Aufg. 3
	% https://mathematikalpha.de/wp-content/uploads/2016/05/aufg125.pdf
	Bestimmen Sie die Werte der folgenden Integrale, vorausgesetzt sie existieren:
	\begin{eqnarray*}
		\hspace{-0.3cm}
		\begin{array}{l l l l}
			(a) \; \displaystyle{\int_0^4} \left( x^3 - 5x^2 + \frac32 x -10\right)\, \mathrm{d}x & (b)\; \displaystyle{\int_1^4} \displaystyle{\frac{1-z^2}{z}}\,\mathrm{d}z &(c)\; \displaystyle{\int_0^{0,5}} \displaystyle{\frac{3}{\sqrt{1-x^2}}}\, \mathrm{d}x & (d)\; \displaystyle{\int_0^{\pi/4}} \displaystyle{\frac{1-\cos^2(x)}{2\cos^2(x)}}\,\mathrm{d}x \\[2ex]
			(e) \; \displaystyle{\int_1^\infty} \frac{1}{x^4}\,\mathrm{d} x & (f)\; \displaystyle{\int_0^\infty} \displaystyle{\frac{1}{x}}\,\mathrm{d}x  & (g)\; \int\limits_{-\infty}^0 e^{-t}\,\mathrm{d}t &(h)\; \displaystyle{\int_0^{4}} \displaystyle{\frac{2}{\sqrt{t}}}\, \mathrm{d}t
		\end{array}
	\end{eqnarray*}

\vspace{0.3cm}
\textbf{L�sung:} \newline
Die hier gestellten Aufgabenteile \( (a)- (d) \) finden Sie in Papula, Band 1, V.1-7, ab Seite \(560\) und die L�sungen auf Seite \(787\)
\begin{itemize}
	\item[(a)] Papula Aufg. 3.a) $\quad -212/3 $
	\item[(b)] Papula Aufg. 3.c) $\quad \ln(4) - 15/2 $
	\item[(c)] Papula Aufg. 3.i) $\quad \pi/2 $
	\item[(d)] Papula Aufg. 3.j) $\quad (1-\pi/4)/2 $
	\item[(e)]
	\[ \int\limits_1^a \frac{1}{x^4}\,\mathrm{d}x = \Big[ -\frac{1}{3x^3}\Big]_1^a = -\frac{1}{3a^3} + \frac13 \Rightarrow \lim\limits_{a\to\infty} \left( -\frac{1}{3a^3} + \frac13\right) = \frac13\,. \]
	\item[(f)]
		\[ \int \frac{1}{x}\,\mathrm{d}x = \ln(x)+C \Rightarrow \lim\limits_{a \to 0}\int\limits_{a}^{1} \frac1x\,\mathrm{d}x + \lim\limits_{b\to\infty} \int\limits_{1}^{b} \frac1x\,\mathrm{d}x\]
Keines der beiden uneigentlichen Integrale existiert: $\lim\limits_{a\to 0} \ln(a) = -\infty$ und  $\lim\limits_{b\to\infty} \ln(b) = \infty$
	\item[(g)]
		\[ \int\limits_a^0 e^{-t}\,\mathrm{d}t = \Big[ - e^{-t}\Big]_a^0 = -1 + e^{-a} \Rightarrow \lim\limits_{a\to-\infty} \left( -1 + e^{-a}\right) = \lim\limits_{a\to \infty} \left( -1 + e^a\right)  \textnormal{ existiert nicht}\,. \]
	\item[(h)]
\[ \int\limits_a^4 \frac{2}{\sqrt{t}}\,\mathrm{d}t = \Big[ 4 \sqrt{t}\Big]_a^4 = 8 - 4\sqrt{a} \Rightarrow \lim\limits_{a\to 0} \left( 8 - 4\sqrt{a}\right) = 8\,. \]
\end{itemize}
\end{exercisebox}

\begin{exercisebox}[Numerische Integration]
	Das folgende bestimmte Integral soll n�herungsweise bestimmt werden:
	\[
		\int_1^3 x^3 \, \mathrm{d}x = 20\,.
	\]
	\begin{itemize}
		\item[(a)] Bestimmen Sie eine Zerlegung des Intervalls $[1,3]$ in $4$ gleich gro�e Intervalle.
		\item[(b)] Wenden Sie eine Rechteckregel an.
		\item[(c)] Wenden Sie die Trapezregel an.
		\item[(d)] Wenden Sie die Simpson-Formel an.
	\end{itemize}
\vspace{0.3cm}
\textbf{L�sung:}
\begin{itemize}
	\item[(a)] $x_0=1, x_1 = \frac32, x_2 = 2, x_3=\frac52, x_4=3$
	\item[(b)] $A = \frac12\cdot (1^3 + (\frac32)^3 + 2^3 + (\frac52)^3) = 14$
	\item[(c)] $A = \frac14\cdot 1^3 + \frac12( (\frac32)^3 + 2^3)  + \frac14 (\frac52)^3) = 20.5$
\item[(d)] Zur Anwendung der Simpson-Formel brauchen wir noch alle Zwischenstellen: \[ x_{0}=0, x_{01}=\frac54, x_1 = \frac32, x_{12}=\frac74, x_2 = 2, x_{23}=\frac94, x_3=\frac52, x_{34} = \frac{11}{4}, x_4=3\,.\]
	Die Simpson-Formel ergibt: \begin{eqnarray*} A&=&\frac{1}{12} \left( 1^3 + 4 \left( \frac54\right)^3 + 2 \left( \frac32\right)^3 + 4 \left( \frac74\right)^3 + 2 \left( 2\right)^3 + 4  \left( \frac94\right)^3 + 2 \left( \frac52\right)^3 + 4  \left( \frac{11}{4}\right)^3 + 3^3\right) \\
	&=& \cdots = 20\,.\end{eqnarray*}
	Das Ergebnis h�tten wir auch ohne Rechnung sofort angeben k�nnen, denn die Simpson-Formel ist exakt f�r Polynome dritten Grades. Dieses Ergebnis l��t sich sofort der Fehler-Absch�tzung ablesen - warum?
\end{itemize}
\end{exercisebox}

\begin{exercisebox}[Simpson-Regel]
	Bestimmen Sie den relativen Fehler, den Sie bei einer Approximation des Intgrals \[ I = \int\limits_0^1 f(x) \, \mathrm d x = \int\limits_0^1 x\cdot \cos(x^2) \, \mathrm d x\] mit Hilfe der Simpson-Regel auf drei �quidistanten Teilintervallen machen.

\vspace{0.3cm}
\textbf{L�sung:}
Auszuwerten ist die Simpsonregel f�r \(h=\dfrac13\). Zus�tzlich zu den St�tzpunkten in Intervallr�ndern kommen St�tzpunkte in jeder Intervallmitte. Insgesamt sind das also \(7\) St�tzpunkte \[ \begin{cases} x_k &= 0+ k\cdot \dfrac12 h = k\cdot \dfrac16 \,, \\ y_k &= f(x_k) = x_k \cdot \cos\left( x_k^2\right) \,. \end{cases} \] Mit dieser Notation lautet die Simpsonregel also
\[ S_3(f) = h\left( \dfrac{1}{6} y_0 +  \dfrac23 y_1 + \dfrac13 y_2 + \dfrac23 y_3 + \dfrac13 y_4 + \dfrac23 y_5 + \dfrac16 y_6 \right) = \sum\limits_{k=0}^6 w_k y_k, \]
\small{
\begin{tabular}{c | c c c c c c c}
\(k\) & \(0\) & \(1\) & \(2\) & \(3\) & \(4\) & \(5\) & \(6\) \\\hline
\(x_k\) & \(0\) & \(\dfrac16\) & \(\dfrac26\) & \(\dfrac36\) & \(\dfrac46\) & \(\dfrac56\) & \(\dfrac66\) \\[2ex]
\(y_k\) & \(\dfrac16\cos\left(\dfrac{1}{36}\right)\) & \(\dfrac13\cos\left(\dfrac{1}{9}\right)\) & \(\dfrac12\cos\left(\dfrac14\right)\) & \(\dfrac23\cos\left( \dfrac49\right)\) & \(\dfrac23\cos\left( \dfrac49\right)\) & \(\dfrac56\cos\left(\dfrac{25}{36}\right)\) & \(\cos\left(1\right)\) \\[2ex]
\(w_k\) & \(\dfrac13 \cdot \dfrac16\) & \(\dfrac13 \cdot \dfrac23\)& \(\dfrac13 \cdot \dfrac13\) & \(\dfrac13 \cdot \dfrac23\) & \(\dfrac13 \cdot \dfrac13\) & \(\dfrac13 \cdot \dfrac23\)  & \(\dfrac13 \cdot \dfrac16\)
\end{tabular}
}

Damit \[ S_{3}(f) = \sum\limits_{k=0}^6 w_k \cdot y_k = 0.420681 \,.\]
Analytische Integration liefert \[ I = \int\limits_0^1 x\cdot \cos(x^2)\,\mathrm{d}x = \Big[ \dfrac12 \sin(x^2)\Big]_0^1 = 0.420736\,.\]
Damit ergibt sich f�r den relativen Fehler:
\[e= \dfrac{ | I - S_{3} |}{ I} = 1.31\cdot 10^{-4}\,.\]
\end{exercisebox}

